%
% methode.tex -- Methode der Integralberechnung
%
% (c) 2021 Prof Dr Andreas Müller, OST Ostschweizer Fachhochschule
%
\documentclass[tikz]{standalone}
\usepackage{amsmath}
\usepackage{times}
\usepackage{txfonts}
\usepackage{pgfplots}
\usepackage{csvsimple}
\usetikzlibrary{arrows,intersections,math,calc}
\definecolor{darkred}{rgb}{0.8,0,0}
\begin{document}
\def\skala{1}
\def\kreuz#1{
	\draw[color=blue,line width=1.0pt] ($#1+(-0.1,-0.1)$) -- ++(0.2,0.2);
	\draw[color=blue,line width=1.0pt] ($#1+(-0.1,0.1)$) -- ++(0.2,-0.2);
}
\begin{tikzpicture}[>=latex,thick,scale=\skala,
declare function = {
	sqr(\t) = (\t) * (\t);
	X(\t) = 4 - 2.5 * cos(180*\t) + 7*(\t-0.5)*sin(180*\t);
	Y(\t) = 2.5*sin(180*\t) + \t * (1 - \t) * sqr(1.2 * cos(2*360*\t) + 0.6*sin(3*360*\t) + 1.0*cos(4*360*\t));
}]

\fill[color=darkred!10,line width=1.4pt] plot[domain=0:1,samples=200]
	({X(\x)},{Y(\x)})
	-- cycle;

\coordinate (A) at (1.5,0);
\coordinate (B) at (6.5,0);

\fill[color=gray!20] (A) rectangle ($(B)+(0,-1.1)$);

\draw[->] (-0.1,0) -- (8.4,0) coordinate[label={$\operatorname{Re}$}];
\draw[->] (0,-0.1) -- (0,4.4) coordinate[label={right:$\operatorname{Im}$}];


\draw[color=darkred,line width=1.4pt] plot[domain=0:1,samples=200]
	({X(\x)},{Y(\x)});

\draw[color=darkred,line width=1.4pt] (A) -- (B);
\fill[color=darkred] (A) circle[radius=0.08];
\fill[color=darkred] (B) circle[radius=0.08];
\node at (A) [below left] {$a\mathstrut$};
\node at (B) [below right] {$b\mathstrut$};
\node at ($0.5*(A)+0.5*(B)$) [below] {$\displaystyle\int_a^b f(x)\,dx$};

\coordinate (Z0) at (2,1.1);
\coordinate (Z1) at (3.9,2.5);
\coordinate (Z2) at (5.5,1.3);

\kreuz{(Z0)}
\kreuz{(Z1)}
\kreuz{(Z2)}

\node[color=blue] at (Z0) [left] {$z_0$};
\node[color=blue] at (Z1) [left] {$z_1$};
\node[color=blue] at (Z2) [right] {$z_2$};

\def\r{0.6}

\draw[color=blue] (Z0) circle[radius=\r];
\draw[color=blue] (Z1) circle[radius=\r];
\draw[color=blue] (Z2) circle[radius=\r];

\draw[<-,color=blue] ($(Z0)+(45:{0.99*\r})+(135:0.1)$) -- ++(-45:0.01);
\draw[<-,color=blue] ($(Z1)+(45:{0.99*\r})+(135:0.1)$) -- ++(-45:0.01);
\draw[<-,color=blue] ($(Z2)+(45:{0.99*\r})+(135:0.1)$) -- ++(-45:0.01);

\def\t{0.49}
\node[color=darkred] at ({X(\t)},{Y(\t)}) [above] {$\gamma_0$};

\draw[->,color=darkred,line width=1.2pt]
	($0.3333*(A)+0.6666*(B)$) -- ++(0.1,0);
\draw[->,color=darkred,line width=1.2pt]
	($0.6666*(A)+0.3333*(B)$) -- ++(0.1,0);

\def\t{0.87}
\draw[->,color=darkred,line width=1.2pt]
	({X(\t)},{Y(\t)})
	--
	({X(\t-0.001)},{Y(\t-0.001)});

\def\t{0.68}
\draw[->,color=darkred,line width=1.2pt]
	({X(\t)},{Y(\t)})
	--
	({X(\t-0.001)},{Y(\t-0.001)});

\def\t{0.53}
\draw[->,color=darkred,line width=1.2pt]
	({X(\t)},{Y(\t)})
	--
	({X(\t-0.001)},{Y(\t-0.001)});

\def\t{0.43}
\draw[->,color=darkred,line width=1.2pt]
	({X(\t)},{Y(\t)})
	--
	({X(\t-0.001)},{Y(\t-0.001)});

\def\t{0.24}
\draw[->,color=darkred,line width=1.2pt]
	({X(\t)},{Y(\t)})
	--
	({X(\t-0.001)},{Y(\t-0.001)});

\def\t{0.13}
\draw[->,color=darkred,line width=1.2pt]
	({X(\t)},{Y(\t)})
	--
	({X(\t-0.001)},{Y(\t-0.001)});

\end{tikzpicture}
\end{document}

